\documentclass[10pt,a4paper]{amsart}
\usepackage[margin=19mm]{geometry}
\usepackage{amsmath,amssymb,mathtools}
\usepackage[scr=rsfs,cal=euler]{mathalfa}
\usepackage{xfrac}
\usepackage[unicode,hidelinks,bookmarks=false]{hyperref}
\newcommand{\ie}{i.e.\ }
\newcommand{\cf}{cf.\ }
\newcommand{\resp}{respectively\ }
\newcommand{\Subsection}{\subsection}
\providecommand{\nobreakdash}{\nobreak\hbox{-}}
\setlength{\emergencystretch}{2em}
\multlinegap0pt
\def\R{{\mathbb R}}
\theoremstyle{plain}
\newtheorem{theorem}{Theorem}[section]
\newtheorem{lemma}[theorem]{Lemma}
\newtheorem{proposition}[theorem]{Proposition}
\newtheorem{corollary}[theorem]{Corollary}
\newtheorem{definition}[theorem]{Definition}
\newtheorem{remark}[theorem]{Remark}
\newtheorem{conjecture}[theorem]{Conjecture}
\title[Nonlinear Landau damping for the unconfined Vlasov-Poisson system in 2D]{Nonlinear Landau damping for the unconfined Vlasov-Poisson system in 2D}
\author{Alexandru D. Ionescu, Quoc-Hung Nguyen}
\date{Source: arXiv:2609.28251v1 (CC BY 4.0)}
\begin{document}
\maketitle

\noindent\emph{Source and licence.} The delivered work is the article shown in the title above, version arXiv:2609.28251v1 (CC BY 4.0), distributed under the Creative Commons Attribution 4.0 International licence, as recorded in the arXiv licence metadata for this exact version and shown on the abstract page saved with this item. The full licence notice, which carries the licence URL and the address of that abstract page, is the bundled file \texttt{licenses/arxiv-2609.28251-CC-BY-4.0.txt}.
\par\smallskip

\noindent\textit{Editorial scope.} Counted unit: the article's Lemma (source label keyle) (source0.tex lines 1472--1493, source label \texttt{keyle}): ``Estimates for the characteristics For all we have Y s,t s sY s,t M T a 0(s,v), xY s,t s x sY s,t M T a 1(s,v), vY s,t s v sY s,t M T ( v -1 s - 1 6 - 0 s, v 3 2 - 1 3 0 ), x vY s,t s x v sY s,t s ( x ''. It is delivered together with the article's complete original proof (source0.tex lines 1496--1869, one proof environment printed later in the article, detached from the statement by the article's own arrangement, 374 lines), copied line for line from the article's own text. Required context retained uncounted: de1 (lines 831--836); bo1 (lines 880--882); a23 (lines 946--954); z216 (lines 1078--1108); Yinteq (lines 1451--1453). Editorial formatting only: an AMS \texttt{amsart} preamble that restates the article's own macro definitions and its own theorem declarations, and the theorem environments the retained lines use; the article's own class file is neither shipped nor input; the retained environments are renumbered sequentially in order of appearance, whereas the article prints them with its own numbering; the article's equation numbering is not reproduced and every cross-reference inside the retained text resolves to the wrapper's numbering. No citation occurs inside any retained line, so the wrapper carries no bibliography; All retained bodies are exact copies of the bundled \texttt{sources/211565/source0.tex}, so no omission receipt applies. No asset-inclusion command occurs inside any retained span and the package ships no image asset.


\section*{Retained context: de1 (source0.tex lines 831--836)}

	\begin{equation}\label{de1}
		\begin{split}
			&E(t)=	E^{reg}(t)+E^{osc}(t)=E^{reg}(t)+\cos(t) E^{\cos}(t)+\sin(t)E^{\sin}(t),\\
			&E^{\cos}:=E^{\cos}_{non}+E_{lin}^{\cos},\qquad E^{\sin}:=E^{\sin}_{non}+E_{lin}^{\sin},
		\end{split}
	\end{equation}


\section*{Retained context: bo1 (source0.tex lines 880--882)}

		\begin{equation}\label{bo1}
			M_T := \|E^{reg}\|_{1,T} + \|(E^{\cos}_{non},E^{\sin}_{non})\|_{2,T} + [f_0] \le \epsilon_0,
		\end{equation}


\section*{Retained context: a23 (source0.tex lines 946--954)}

		\begin{align}\label{a23}
			|\nabla_{x,t}^k E_{lin}^{\ast}(t,x)|
			\lesssim [f_0]
			\begin{cases}
				\langle t,x\rangle^{-2-k},&k\in\{0,1\},\\
				\langle t\rangle^{-k+1+20\kappa_0}
				\langle t,x\rangle^{-3-20\kappa_0},&k\in\{2,3\},
			\end{cases}
		\end{align}


\section*{Retained context: z216 (source0.tex lines 1078--1108)}

	\begin{lemma}\label{z216}
		(i) If $0\le s\le t$ and $x,v\in\R^2$, then
		\begin{equation*}
			\int_s^t \langle\tau,x-\tau v\rangle^{-1}\,d\tau
			\lesssim
			\langle v\rangle^{-1}\log\!\Big(2+\frac{t\langle v\rangle}{s+1}\Big).
		\end{equation*}
		
		(ii) Assume that $a_1>1$, $a_2\in\mathbb R$, and set
			$\gamma:=a_2-a_1+1.$
	 If $\gamma>0$ and $0\leq s\leq t$, then
		\begin{equation}\label{z210}
			\int_s^t\langle\tau\rangle^{a_2}\langle\tau,x-\tau v\rangle^{-a_1}\,d\tau\lesssim\langle t\rangle^{\gamma}
			\langle v\rangle^{-1}.
		\end{equation}
		On the other hand, if $\gamma<0$, then \begin{equation}\label{z212}
			\int_s^t
			\langle\tau\rangle^{a_2}
			\langle\tau,x-\tau v\rangle^{-a_1}\,d\tau
			\lesssim
			\langle s\rangle^{\gamma}
			\langle v\rangle^{-1}.
		\end{equation}
	If $\gamma=0$, then
		\begin{equation}\label{z212-borderline}
			\int_s^t\langle\tau\rangle^{a_2}
			\langle\tau,x-\tau v\rangle^{-a_1}\,d\tau
			\lesssim\langle v\rangle^{-1}
			\log\!\left(2+\frac{t+1}{s+1}\right).
		\end{equation}
	\end{lemma}


\section*{Retained context: Yinteq (source0.tex lines 1451--1453)}

	\begin{equation}\label{Yinteq}
	Y_{s,t}(x,v)=\int_s^t(\tau-s)\,E(\tau,x+\tau v+Y_{\tau,t}(x,v))\,d\tau.
	\end{equation}


\section{The counted unit: Lemma (source label keyle) and its complete original proof}

	\begin{lemma} [Estimates for the characteristics]\label{keyle}  
	For all $0\le s\le t\le T$ we have
	\begin{align}\label{n13}
		&	|Y_{s,t}|+\langle s\rangle	|\partial_sY_{s,t}|\lesssim   M_T\mathfrak{a}_0(s,v),\\
		&\label{n14}	|\nabla_xY_{s,t}|+\langle s\rangle	|\nabla_x\partial_sY_{s,t}|\lesssim  M_T\mathfrak{a}_1(s,v),\\
		&\label{xf}|\nabla_vY_{s,t}|+\langle s\rangle	|\nabla_v\partial_sY_{s,t}|\lesssim  M_T\big(\langle v\rangle^{-1}\langle s\rangle^{-\frac{1}{6}-\kappa_0}+\langle s,|v|^{\frac{3}{2}}\rangle^{-\frac{1}{3}+\kappa_0}\big),\\
		&\label{n15}	|\nabla_x\nabla_vY_{s,t}|+\langle s\rangle 	|\nabla_x\nabla_v\partial_sY_{s,t}|+\langle s\rangle(	|\nabla_x^2Y_{s,t}|+\langle s\rangle	|\nabla_x^2\partial_sY_{s,t}|)\lesssim  M_T\langle s\rangle\mathfrak{a}_2(s,v),\\
		&\label{n17}|\nabla_v^2Y_{s,t}|\lesssim M_T	\langle t\rangle^{\frac{2}{3}+\kappa_0}	\langle t^{\frac{1}{2}+2\kappa_0},v\rangle^{-1}
		,\\&\label{n17b}
		|\nabla_v^2\partial_sY_{s,t}|\lesssim	M_T\langle t\rangle^{\frac{1}{6}-\kappa_0}	\langle s^{\frac12+2\kappa_0},v\rangle^{-1}
		\langle s\rangle^{-\frac12+2\kappa_0}.
	\end{align}
	Here, $Y_{s,t}=Y_{s,t}(x,v)$. In particular, the following
	rougher bounds also hold:
	\begin{align}\label{n13'}
		&	|Y_{s,t}|+\langle s\rangle	|\partial_sY_{s,t}|\lesssim   M_T
		\langle v\rangle^{-\frac{1}{10}}\langle s\rangle^{-\frac12-\kappa_0}	,\\&\label{n14'}	|\nabla_xY_{s,t}|+\langle s\rangle	|\nabla_x\partial_sY_{s,t}|\lesssim  M_T\langle v\rangle^{-\frac{1}{10}}\langle s\rangle^{-\frac{7}{6}-\kappa_0}
		,\\&\label{xf'}
		|\nabla_vY_{s,t}|+\langle s\rangle	|\nabla_v\partial_sY_{s,t}|\lesssim  M_T\langle v\rangle^{-\frac{1}{10}}\langle s\rangle^{-\frac{1}{6}-\kappa_0},\\
		&\label{n15'}	|\nabla_x\nabla_vY_{s,t}|+\langle s\rangle 	|\nabla_x\nabla_v\partial_sY_{s,t}|+\langle s\rangle(	|\nabla_x^2Y_{s,t}|+\langle s\rangle	|\nabla_x^2\partial_sY_{s,t}|)\lesssim  M_T\langle v\rangle^{-2\kappa_0}  \langle s\rangle^{-\frac{5}{6}+2\kappa_0}.
	\end{align}
\end{lemma}

\begin{proof} We divide the proof into three steps.
	\medskip
	
	\textbf{Step 1: oscillatory integral bounds.} For integers $n_1,n_2\ge0$ with $n_1+n_2\le2$, define
	\begin{align*}
		P_{n_1,n_2}	:=	\Big|\int_s^t(\tau-s) \tau^{n_1}\nabla_{x}^{n_1+n_2}E^{osc}(\tau,x+\tau v)d\tau\Big|+\langle s\rangle\Big|\int_s^t \tau^{n_1}\nabla_{x}^{n_1+n_2}E^{osc}(\tau,x+\tau v)\,d\tau\Big|.
	\end{align*}
	Using \eqref{de1} and integrating by parts on $[s_0,t]$ for any $s_0\in[s,t]$, we obtain
	\begin{equation}\label{Pbounds1}
		\begin{split}
			P_{n_1,n_2}	&\lesssim \sum_{\ast\in\{\cos,\sin\}}\Big\{\int_s^{s_0} \langle \tau\rangle^{n_1+1}|\nabla_{x}^{n_1+n_2}E^{\ast}(\tau,x+\tau v)|\,d\tau\\
			&+\mathbf 1_{s_0<t}\sum_{\tau\in\{s_0,t\}} \langle \tau\rangle^{n_1+1}|\nabla_{x}^{n_1+n_2} E^{\ast}(\tau,x+\tau v)|+\int_{s_0}^{t}\langle \tau\rangle^{n_1}\Big[\langle \tau\rangle|\nabla_{x}^{n_1+n_2}\partial_\tau E^{\ast}(\tau,x+\tau v)|\\
			&+|\nabla_{x}^{n_1+n_2} E^{\ast}(\tau,x+\tau v)|+\langle v\rangle\langle \tau\rangle|\nabla_{x}^{n_1+n_2+1} E^{\ast}(\tau,x+\tau v)|\Big]\,d\tau\Big\}.
		\end{split}
	\end{equation}
	
	Using \eqref{bo1} and \eqref{a23}, for any $t\in[0,T]$, $x\in\R^2$, and $\ast\in\{\cos,\sin\}$ we obtain
	\begin{equation}\label{Ebounds1}
		\begin{split}
			&|E^\ast(t,x)|\lesssim M_T\langle t,x\rangle^{-2}\langle t\rangle^{50\kappa_0},\\
			&|\nabla_x^n E^\ast(t,x)|\lesssim M_T\langle t,x\rangle^{-2-\kappa_0}\langle t\rangle^{-2n/3+2\kappa_0},\qquad n\in\{1,2\},\\
			&|\partial_t\nabla_x^n E^\ast(t,x)|\lesssim M_T\langle t,x\rangle^{-2-\kappa_0}\langle t\rangle^{-1/2-2n/3},\qquad n\in\{0,1\},\\
			&|\nabla_x^3 E^\ast(t,x)|+|\partial_t\nabla_x^2 E^\ast(t,x)|\lesssim M_T\langle t,x\rangle^{-2-\kappa_0}\langle t\rangle^{-11/6}.
		\end{split}
	\end{equation}
	
	
	{\bf{Case $n_1+n_2=0$.}} Using \eqref{Pbounds1} and \eqref{Ebounds1}, we have
	\begin{align*}
		P_{0,0}&\lesssim
		M_T\int_s^{s_0}\langle \tau\rangle^{1+50\kappa_0}
		\langle \tau,x+\tau v\rangle^{-2}\,d\tau+M_T\mathbf 1_{s_0<t}\sum_{r\in\{s_0,t\}}
		\langle r\rangle^{1+50\kappa_0}
		\langle r,x+rv\rangle^{-2}\\
		&+M_T\int_{s_0}^{t}\Big[
		\langle \tau\rangle^{\frac12}
		\langle \tau,x+\tau v\rangle^{-2-\kappa_0}+
		\langle \tau\rangle^{50\kappa_0}
		\langle \tau,x+\tau v\rangle^{-2}+
		\langle v\rangle
		\langle \tau\rangle^{\frac13+\kappa_0}
		\langle \tau,x+\tau v\rangle^{-2}
		\Big]\,d\tau .
	\end{align*}
	Applying \eqref{z210}--\eqref{z212}, we obtain
	\begin{align}
		P_{0,0}&\lesssim M_T\big(\mathbf 1_{s_0>s}\langle v\rangle^{-1}
		\langle s_0\rangle^{50\kappa_0}+\mathbf 1_{s_0<t}
		\langle v\rangle^{-1}
		\langle s_0\rangle^{-\frac12-\kappa_0}+\mathbf 1_{s_0<t}
		\langle s_0\rangle^{-\frac23+\kappa_0}\big).
		\label{P00-pre}
	\end{align}
	We choose $s_0\in[s,t]$ as 
	\begin{equation}\label{choice-s0-P00}
		s_0:=\begin{cases}
			s\qquad&\text{ if }\,\,|v|^{3/2}\leq s,\\
			t\qquad&\text{ if }\,\,|v|^{3/2}\geq t,\\
			|v|^{3/2}&\text{ if }\,\,s\leq |v|^{3/2}\leq t.
		\end{cases}		
	\end{equation}
	With this choice,
	\begin{align*}
		\mathbf 1_{s_0>s}
		\langle v\rangle^{-1}
		\langle s_0\rangle^{50\kappa_0}
		&\lesssim
		\langle s,\langle v\rangle^{3/2}\rangle^{-\frac23+50\kappa_0},\\
		\mathbf 1_{s_0<t}\langle s_0\rangle^{-\frac23+\kappa_0}
		&\lesssim
		\langle s,\langle v\rangle^{3/2}\rangle^{-\frac23+\kappa_0}\lesssim
		\langle s,\langle v\rangle^{3/2}\rangle^{-\frac23+50\kappa_0},\\
		\mathbf 1_{s_0<t}
		\langle v\rangle^{-1}
		\langle s_0\rangle^{-\frac12-\kappa_0}
		&\lesssim
		\langle v\rangle^{-1}
		\langle s\rangle^{-\frac12-\kappa_0}.
	\end{align*}
	Therefore
	\begin{equation}\label{q60}
		P_{0,0}\lesssim M_T\big(\langle v\rangle^{-1}\langle s\rangle^{-\frac12-\kappa_0}+
		\langle s,|v|^{3/2}\rangle^{-\frac23+50\kappa_0}\big).
	\end{equation}
	
	
	
	{\bf{Case $n_1+n_2=1$.}} In this case we prove that
	\begin{align}\nonumber
		P_{n_1,n_2}&\lesssim \mathbf{1}_{n_1=0}M_T\big(\langle v\rangle^{-1}\langle s\rangle^{-\frac{7}{6}-\kappa_0}+\langle s,|v|^{\frac{3}{2}}\rangle^{-\frac{2}{3}} 	\langle s\rangle^{-\frac{2}{3}+\kappa_0}\big)
		\\&+\mathbf{1}_{n_1=1}M_T\big(\langle v\rangle^{-1}\langle s\rangle^{-\frac{1}{6}-\kappa_0}+\langle s,|v|^{\frac{3}{2}}\rangle^{-\frac{1}{3}+\kappa_0}\big).\label{q61}
	\end{align}	
	Indeed, from \eqref{Pbounds1} and \eqref{Ebounds1}, we have
	\begin{align}\nonumber
		P_{n_1,n_2}	&\lesssim M_T\int_s^{s_0} \langle \tau\rangle^{\frac{1}{3}+n_1+2\kappa_0} \Phi(\tau,x+\tau v)d\tau+\mathbf 1_{s_0<t}\sum_{\tau\in\{s_0,t\}}M_T \langle \tau\rangle^{\frac{1}{3}+n_1+2\kappa_0} \Phi(\tau,x+\tau v)\\&
		+M_T\int_{s_0}^{t}(\langle \tau\rangle^{n_1-\frac{1}{6}}+\langle v\rangle\langle \tau\rangle^{n_1-\frac{1}{3}+2\kappa_0})\Phi(\tau,x+\tau v)d\tau.	\label{P-one-derivative-start}
	\end{align}
	
	We apply now \eqref{z210}--\eqref{z212}. First, for the integral over $[s,s_0]$, Lemma~\ref{z216} gives
	\begin{align}
		&\int_s^{s_0}
		\langle \tau\rangle^{\frac13+n_1+2\kappa_0}
		\Phi(\tau,x+\tau v)\,d\tau\lesssim\mathbf 1_{s_0>s}\langle v\rangle^{-1}
		\big(\mathbf 1_{n_1=0}\langle s\rangle^{-\frac23+\kappa_0}+\mathbf 1_{n_1=1}\langle s_0\rangle^{\frac13+\kappa_0}\big)
		\label{q61-I0}.
	\end{align}
	Second, for the integrals over $[s_0,t]$, we obtain
	\begin{equation*}
		\begin{split}
			\int_{s_0}^{t}
			\langle \tau\rangle^{n_1-\frac16}
			\Phi(\tau,x+\tau v)\,d\tau
			&\lesssim\mathbf{1}_{s_0<t}\langle v\rangle^{-1}\langle s_0\rangle^{n_1-\frac76-\kappa_0},\\
			\langle v\rangle
			\int_{s_0}^{t}
			\langle \tau\rangle^{n_1-\frac13+2\kappa_0}
			\Phi(\tau,x+\tau v)\,d\tau
			&\lesssim\mathbf{1}_{s_0<t}\langle s_0\rangle^{n_1-\frac43+\kappa_0}.
		\end{split}
	\end{equation*}
	Finally, $\langle r\rangle^{\frac13+n_1+2\kappa_0}\Phi(r,x+rv)\lesssim\langle r\rangle^{n_1-\frac53+\kappa_0}$ for $r\in\{s_0,t\}$. Thus, 
	\begin{align}\nonumber
		P_{n_1,n_2}	&\lesssim M_T	\mathbf 1_{s_0>s}\langle v\rangle^{-1}\big(\mathbf 1_{n_1=0}
		\langle s\rangle^{-\frac23+\kappa_0}+\mathbf 1_{n_1=1}
		\langle s_0\rangle^{\frac13+\kappa_0}\big)\\&
		+M_T\mathbf{1}_{s_0<t}\big(\langle v\rangle^{-1}
		\langle s_0\rangle^{n_1-\frac76-\kappa_0}+	\langle s_0\rangle^{n_1-\frac43+\kappa_0}\big).	\label{P-one-derivative-start'}
	\end{align}
	We now choose  $s_0$ as in \eqref{choice-s0-P00} and the desired bounds \eqref{q61} follow from \eqref{P-one-derivative-start'}.
	
	{\bf{Case $n_1+n_2=2$.}} In this case we prove that
	\begin{align}
		P_{n_1,n_2}&\lesssim\mathbf{1}_{n_1\leq 1}M_T
		\langle s^{\frac12+2\kappa_0},v\rangle^{-1}
		\langle s\rangle^{-\frac43+n_1+\kappa_0}+\mathbf{1}_{n_1=2}M_T
		\langle t\rangle^{\frac23+\kappa_0}
		\langle v,t^{\frac12+2\kappa_0}\rangle^{-1}
		\label{q62}.
	\end{align}
	From \eqref{Pbounds1} and \eqref{Ebounds1}, we have, for $n_1+n_2=2$,
	\begin{equation}\label{P-two-derivative-start}
		\begin{split}
			&P_{n_1,n_2}\lesssim
			M_T\int_s^{s_0}
			\langle \tau\rangle^{-\frac13+n_1+2\kappa_0}\Phi(\tau,x+\tau v)\,d\tau\\
			&+\mathbf{1}_{s_0<t}M_T\sum_{r\in\{s_0,t\}}
			\langle r\rangle^{-\frac13+n_1+2\kappa_0}\Phi(r,x+rv)+M_T\int_{s_0}^{t}
			\langle v\rangle
			\langle \tau\rangle^{n_1-\frac56}
			\Phi(\tau,x+\tau v)\,d\tau.
		\end{split}	
	\end{equation}
	
	Assume first that $n_1\le1$. Applying \eqref{z210}--\eqref{z212} gives
	\begin{align}
		P_{n_1,n_2}
		&\lesssim
		M_T\mathbf 1_{s_0>s}
		\langle v\rangle^{-1}
		\langle s\rangle^{n_1-\frac43+\kappa_0}+
		M_T\mathbf 1_{s_0<t}
		\langle s_0\rangle^{n_1-\frac{11}{6}-\kappa_0}.
		\label{q62-reduced-low}
	\end{align}
	Choose $s_0\in[s,t]$ such that
	\begin{equation*}
		s_0=t
		\quad\text{if}\quad
		\langle v\rangle\ge \langle s\rangle^{\frac12+2\kappa_0},
		\qquad
		s_0=s
		\quad\text{if}\quad
		\langle v\rangle< \langle s\rangle^{\frac12+2\kappa_0}.
	\end{equation*}
	The bounds \eqref{q62} follow in this case from \eqref{q62-reduced-low}.
	
	
	It remains to consider $n_1=2$, $n_2=0$. Applying
	\eqref{z210} to \eqref{P-two-derivative-start}, we obtain
	\begin{align*}
		P_{2,0}&\lesssim M_T\mathbf 1_{s_0>s}\langle v\rangle^{-1}
		\langle s_0\rangle^{\frac23+\kappa_0}+M_T\mathbf 1_{s_0<t}
		\big(\langle t\rangle^{\frac16-\kappa_0}+\langle s_0\rangle^{-\frac13+\kappa_0}\big).
	\end{align*}
	We choose $s_0\in[s,t]$ such that
	\[
	s_0=t
	\quad\text{if}\quad
	\langle v\rangle\geq
	\langle t\rangle^{\frac12+2\kappa_0},
	\qquad
	s_0=s
	\quad\text{if}\quad
	\langle v\rangle<
	\langle t\rangle^{\frac12+2\kappa_0},
	\]
	and the desired bounds \eqref{q62} follow in this case as well.
	
	For later use we notice that our argument also proves that
	\begin{align}
		\left|\int_s^t\tau^2\nabla_x^2E^{osc}(\tau,x+\tau v)\,d\tau\right|
		&\lesssim M_T\langle s^{\frac12+2\kappa_0},v\rangle^{-1}\langle s\rangle^{-\frac13+\kappa_0}.
		\label{q62p}
	\end{align}
	Indeed, $\tau^2=(\tau-s)\tau+s\tau$, so the term in the left-hand side of \eqref{q62p} is controlled by $P_{1,1}$.
	\medskip
	
	\textbf{Step 2: bootstrap control.}
	Assume that for some $T'\le T$,
	\begin{equation}\label{n12}
		\begin{split}
			&A(T'):=	\sup_{0\leq s\leq t\leq T'}\big\{\langle s\rangle^{\frac{1}{2}+\kappa_0}\|(Y_{s,t},\langle s\rangle\partial_sY_{s,t})\|_{L^\infty_{x,v}}+\langle s\rangle^{\frac{7}{6}+\kappa_0}\|(\nabla_xY_{s,t},\langle s\rangle\partial_s\nabla_xY_{s,t})\|_{L^\infty_{x,v}}\\
			&+\langle s\rangle^{\frac{1}{6}+\kappa_0}\|(\nabla_vY_{s,t},\langle s\rangle\partial_s\nabla_vY_{s,t})\|_{L^\infty_{x,v}}	+\langle s\rangle^{\frac{5}{6}+\kappa_0}\|(\nabla_x\nabla_vY_{s,t},\langle s\rangle\partial_s\nabla_x\nabla_vY_{s,t})\|_{L^\infty_{x,v}}\\
			&+\langle s\rangle^{\frac{11}{6}+\kappa_0}\|(\nabla_x^2Y_{s,t},\langle s\rangle\partial_s\nabla_x^2Y_{s,t})\|_{L^\infty_{x,v}}+\langle t\rangle^{-\frac{1}{6}+\kappa_0}\|(\nabla_v^2Y_{s,t},\langle s\rangle\partial_s\nabla_v^2Y_{s,t})\|_{L^\infty_{x,v}}\big\}\leq 1.
		\end{split}
	\end{equation}
	
	For simplicity of notation we write $\nabla_x^n\widetilde{E}(\tau):=(\nabla_x^nE)(\tau,x+\tau v+Y_{\tau,t})$, $\nabla_x^n\widetilde{E}^{reg}(\tau):=(\nabla_x^nE^{reg})(\tau,x+\tau v+Y_{\tau,t})$, $\nabla_x^n\widetilde{E}^{osc}(\tau):=(\nabla_x^nE^{osc})(\tau,x+\tau v+Y_{\tau,t})$. Differentiating the characteristic equation \eqref{Yinteq} yields the integral identities
	\begin{equation}\label{nabla1Y}
		\begin{split}
			\nabla_xY_{s,t}&=\int_s^t(\tau-s)(I+\nabla_{x}Y_{\tau,t})\nabla \widetilde{E}(\tau)d\tau,\\
			\nabla_vY_{s,t}&=\int_s^t(\tau-s)(\tau I+\nabla_{v}Y_{\tau,t})\nabla \widetilde{E}(\tau)d\tau,
		\end{split}
	\end{equation}
	\begin{equation}\label{nabla2Y}
		\begin{split}
			\nabla_x^2Y_{s,t}&=\int_s^t(\tau-s)\nabla_{x}^2Y_{\tau,t}\nabla \widetilde{E}(\tau)d\tau+\int_s^t(\tau-s)(I+\nabla_{x}Y_{\tau,t})^{\otimes 2}:\nabla^2 \widetilde{E}(\tau)d\tau,\\
			\nabla_v^2Y_{s,t}&=\int_s^t(\tau-s)\nabla_{v}^2Y_{\tau,t}\nabla \widetilde{E}(\tau)\,d\tau+\int_s^t(\tau-s)(\tau I+\nabla_{v}Y_{\tau,t})^{\otimes 2}:\nabla^2 \widetilde{E}(\tau)d\tau,\\
			\nabla_x\nabla_{v}Y_{s,t}&=\int_s^t(\tau-s)\nabla_{x}\nabla_vY_{\tau,t}\nabla \widetilde{E}(\tau)d\tau\\
			&+\int_s^t(\tau-s)(I+\nabla_{x}Y_{\tau,t})\otimes (\tau I+\nabla_{v}Y_{\tau,t}):\nabla^2 \widetilde{E}(\tau)d\tau.
		\end{split}
	\end{equation}
	
	Using the bootstrap bounds \eqref{n12} together with \eqref{bo1}, we reduce each quantity to an oscillatory part controlled by $P_{n_1,n_2}$ plus a regular contribution estimated by \eqref{z210}--\eqref{z212}. Notice that $|Y_{\tau,t}|\lesssim1$, due to \eqref{n12}. Therefore, using \eqref{bo1}, \eqref{a23}, and \eqref{Ebounds1}, we have
	\begin{equation}\label{Ecomp}
		\begin{split}
			|\nabla_x^q&E^{osc}(\tau,x+\tau v+Y_{\tau,t}(x,v))
			-\nabla_x^qE^{osc}(\tau,x+\tau v)|\\
			&+|\nabla_x^qE^{reg}(\tau,x+\tau v+Y_{\tau,t}(x,v))|\lesssim M_T\langle\tau\rangle^{-\frac12-\frac{2q}{3}}\Phi(\tau,x+\tau v),
		\end{split}
	\end{equation}
	for $q=0,1,2$. 
	
	We would like to use the identities \eqref{Yinteq} and \eqref{nabla1Y}--\eqref{nabla2Y} to close the bootstrap argument. The terms containing only the free factors $I$ and $\tau I$ can be bounded by the corresponding $P_{n_1,n_2}$ (and by
	\eqref{q62p} in the last
	estimate).  Every other term is covered by the displayed bound or contains
	a derivative of $Y_{\tau,t}$ and is bounded by \eqref{n12} and
	\eqref{Ebounds1}; the
	$\nabla_v^2Y_{\tau,t}$ terms are treated separately in the last two bounds. 
	
	More precisely, starting from \eqref{Yinteq} and using \eqref{Ecomp} we can estimate
	\begin{equation}\label{Ysu1}
		\begin{split}
			|Y_{s,t}|+\langle s\rangle	|\partial_sY_{s,t}|&\lesssim \left|\int_s^t(\tau-s) \widetilde{E}^{osc}(\tau)d\tau\right|+\langle s\rangle\left|\int_s^t \widetilde{E}^{osc}(\tau)d\tau\right|+\int_s^t\langle \tau\rangle| \widetilde{E}^{reg}(\tau)|\,d\tau\\&
			\lesssim 	P_{0,0}+M_T\int_s^t\langle \tau\rangle^{\frac{1}{2}}\Phi(\tau,x+\tau v)d\tau.
		\end{split}
	\end{equation}
	
	Similarly, using \eqref{nabla1Y}, \eqref{Ecomp}, and \eqref{Ebounds1} we estimate
	\begin{equation}\label{Ysu2}
		\begin{split}
			|\nabla_xY_{s,t}|+\langle s\rangle|\nabla_x\partial_sY_{s,t}|&\lesssim \Big|\int_s^t(\tau-s)\nabla \widetilde{E}^{osc}(\tau)d\tau\Big|+\langle s\rangle \Big|\int_s^t\nabla \widetilde{E}^{osc}(\tau)d\tau\Big|\\
			&+\int_s^t\langle \tau\rangle|\nabla \widetilde{E}^{reg}(\tau)|\,d\tau+\int_s^t\langle \tau\rangle^{1-\frac{7}{6}-\kappa_0}|\nabla \widetilde{E}(\tau)|\,d\tau\\
			&\lesssim	P_{0,1}+M_T\int_s^t\langle \tau\rangle^{-\frac{1}{6}}\Phi(\tau,x+\tau v)d\tau,
		\end{split}
	\end{equation}
	and
	\begin{equation}\label{Ysu3}
		\begin{split}
			|\nabla_vY_{s,t}|+\langle s\rangle|\nabla_v\partial_sY_{s,t}|&\lesssim \Big|\int_s^t(\tau-s)\tau\nabla \widetilde{E}^{osc}(\tau)d\tau\Big|+\langle s\rangle \Big|\int_s^t \tau\nabla \widetilde{E}^{osc}(\tau)d\tau\Big|\\
			&+\int_s^t\langle \tau\rangle^2|\nabla \widetilde{E}^{reg}(\tau)|\,d\tau+\int_s^t\langle \tau\rangle^{1-\frac{1}{6}-\kappa_0}|\nabla \widetilde{E}(\tau)|\,d\tau\\
			&\lesssim P_{1,0}+M_T\int_s^t\langle \tau\rangle^{\frac{5}{6}}\Phi(\tau,x+\tau v)d\tau.
		\end{split}
	\end{equation}
	
	Similarly, using \eqref{nabla2Y},
	\eqref{Ecomp}, and \eqref{Ebounds1}, we obtain
	\begin{equation}\label{Ysu4}
		\begin{split}
			&|\nabla_x^2Y_{s,t}|+\langle s\rangle|\nabla_x^2\partial_sY_{s,t}|
			\lesssim\Big|\int_s^t(\tau-s)\nabla^2\widetilde E^{osc}(\tau)\,d\tau\Big|
			+\langle s\rangle\Big|\int_s^t\nabla^2\widetilde E^{osc}(\tau)\,d\tau\Big|\\
			&+\int_s^t\langle\tau\rangle|\nabla^2\widetilde E^{reg}(\tau)|\,d\tau+\int_s^t\langle\tau\rangle^{1-\frac{11}{6}-\kappa_0}
			|\nabla\widetilde E(\tau)|\,d\tau+\int_s^t\langle\tau\rangle^{1-\frac{7}{6}-\kappa_0}|\nabla^2\widetilde E(\tau)|\,d\tau\\
			&\lesssim
			P_{0,2}
			+M_T\int_s^t
			\langle\tau\rangle^{-\frac56}
			\Phi(\tau,x+\tau v)\,d\tau,
		\end{split}
	\end{equation}
	and
	\begin{equation}\label{Ysu5}
		\begin{split}
			&|\nabla_x\nabla_vY_{s,t}|+\langle s\rangle|\nabla_x\nabla_v\partial_sY_{s,t}|
			\lesssim\Big|\int_s^t(\tau-s)\tau
			\nabla^2\widetilde E^{osc}(\tau)\,d\tau\Big|+\langle s\rangle
			\Big|\int_s^t\tau\nabla^2\widetilde E^{osc}(\tau)\,d\tau\Big|\\
			&+\int_s^t\langle\tau\rangle^2
			|\nabla^2\widetilde E^{reg}(\tau)|\,d\tau+\int_s^t\langle\tau\rangle^{1-\frac56-\kappa_0}
			|\nabla\widetilde E(\tau)|\,d\tau+\int_s^t
			\langle\tau\rangle^{1-\frac16-\kappa_0}
			|\nabla^2\widetilde E(\tau)|\,d\tau\\
			&\lesssim P_{1,1}+M_T\int_s^t\langle\tau\rangle^{\frac16}\Phi(\tau,x+\tau v)\,d\tau.
		\end{split}
	\end{equation}
	
	To bound the two $v$-derivatives, we notice that $|\nabla_v^2Y_{\tau,t}|\lesssim \langle t\rangle^{\frac16-\kappa_0}$ (due to \eqref{n12}) thus
	\begin{equation}\label{Ysu6}
		\begin{split}
			&|\nabla_v^2Y_{s,t}|\lesssim\Big|\int_s^t(\tau-s)\tau^2\nabla^2\widetilde E^{osc}(\tau)\,d\tau\Big|
			+\int_s^t\langle\tau\rangle^3|\nabla^2\widetilde E^{reg}(\tau)|\,d\tau\\
			&+\int_s^t\langle\tau\rangle^{2-\frac16-\kappa_0}|\nabla^2\widetilde E(\tau)|\,d\tau+
			\langle t\rangle^{\frac16-\kappa_0}\int_s^t\langle\tau\rangle|\nabla\widetilde E(\tau)|\,d\tau\\
			&\lesssim
			P_{2,0}+M_T\int_s^t
			\langle\tau\rangle^{\frac76}
			\Phi(\tau,x+\tau v)\,d\tau+	M_T\langle t\rangle^{\frac16-\kappa_0}
			\int_s^t
			\langle\tau\rangle^{\frac13+2\kappa_0}
			\Phi(\tau,x+\tau v)\,d\tau
		\end{split}
	\end{equation}
	and
	\begin{equation}\label{Ysu7}
		\begin{split}
			&|\nabla_v^2\partial_sY_{s,t}|\lesssim\Big|\int_s^t\tau^2\nabla^2\widetilde E^{osc}(\tau)\,d\tau\Big|
			+\int_s^t\langle\tau\rangle^2|\nabla^2\widetilde E^{reg}(\tau)|\,d\tau\\
			&+\int_s^t\langle\tau\rangle^{1-\frac16-\kappa_0}|\nabla^2\widetilde E(\tau)|\,d\tau+
			\langle t\rangle^{\frac16-\kappa_0}
			\int_s^t|\nabla\widetilde E(\tau)|\,d\tau\lesssim\Big|\int_s^t\tau^2\nabla^2E^{osc}(\tau,x+\tau v)\,d\tau\Big|\\
			&+M_T\int_s^t\langle\tau\rangle^{\frac16}\Phi(\tau,x+\tau v)\,d\tau+M_T\langle t\rangle^{\frac16-\kappa_0}\int_s^t\langle\tau\rangle^{-\frac23+2\kappa_0}\Phi(\tau,x+\tau v)\,d\tau.
		\end{split}
	\end{equation}
	
	
	\textbf{Step 3: conclusion.}
	For any $\beta\neq1+\kappa_0$, Lemma~\ref{z216} (ii) gives
	\begin{equation}\label{IntBouRep}
		\int_s^t
		\langle\tau\rangle^\beta
		\Phi(\tau,x+\tau v)\,d\tau
		\lesssim
		\langle v\rangle^{-1}
		\begin{cases}
			\langle s\rangle^{\beta-1-\kappa_0},
			&\beta<1+\kappa_0,\\
			\langle t\rangle^{\beta-1-\kappa_0},
			&\beta>1+\kappa_0.
		\end{cases}
	\end{equation}
	We apply this with $\beta=1/2, \beta=-1/6, \beta=5/6, \beta=-5/6, \beta=1/6$, and use the bounds \eqref{Ysu1}--\eqref{Ysu5}, and the bounds \eqref{q60}, \eqref{q61}, and \eqref{q62} (with $n_1\leq 1$). The desired bounds \eqref{n13}, \eqref{n14}, \eqref{xf}, and \eqref{n15} follow. Moreover, the bounds \eqref{Ysu6}, \eqref{q62} (with $n_1=2$), and \eqref{IntBouRep} show that
	\begin{align*}
		|\nabla_v^2Y_{s,t}|&\lesssim
		M_T\langle t\rangle^{\frac23+\kappa_0}
		\langle v,t^{\frac12+2\kappa_0}\rangle^{-1}+M_T\langle v\rangle^{-1}
		\langle t\rangle^{\frac16-\kappa_0}\lesssim
		M_T\langle t\rangle^{\frac23+\kappa_0}
		\langle v,t^{\frac12+2\kappa_0}\rangle^{-1},
	\end{align*}
	which proves \eqref{n17}. Finally, the bounds \eqref{Ysu7}, \eqref{q62p}, and \eqref{IntBouRep} show that
	\begin{align*}
		|\nabla_v^2\partial_sY_{s,t}|&\lesssim M_T\langle s^{\frac12+2\kappa_0},v\rangle^{-1}
		\langle s\rangle^{-\frac13+\kappa_0}+M_T\langle v\rangle^{-1}
		\langle s\rangle^{-\frac56-\kappa_0}+ M_T\langle t\rangle^{\frac16-\kappa_0}\langle v\rangle^{-1}\langle s\rangle^{-\frac53+\kappa_0}\\
		&\lesssim
		M_T\langle t\rangle^{\frac16-\kappa_0}
		\langle s^{\frac12+2\kappa_0},v\rangle^{-1}
		\langle s\rangle^{-\frac12+2\kappa_0},
	\end{align*}
	which proves \eqref{n17b}.
	
	It remains to close the bootstrap. The bounds \eqref{n13}--\eqref{n17b} and the definition \eqref{n12} show that \(A(T')\leq C M_T\) whenever \(A(T')\leq1\). Since $A(0)=0$ and $A(T')$ is continuous it follows that $A(T)\lesssim M_T$ provided that $M_T$ is sufficiently small. This completes the proof.
\end{proof}

\end{document}
